\documentclass[a4paper,12pt,twoside,openright,titlepage]{book}

%Additional packages
\usepackage[utf8]{inputenc}
\usepackage[T1]{fontenc}
\usepackage[dutch,english]{babel}
\usepackage{imakeidx}
\usepackage{syntonly}
\usepackage[official]{eurosym}
%\usepackage[graphicx]
\usepackage{graphicx}
\graphicspath{ {./images/} }
\usepackage{float}
\usepackage{xurl}
\usepackage{hyperref}
\hypersetup{colorlinks=true, linkcolor=blue, citecolor=blue, filecolor=blue, urlcolor=blue, pdftitle=, pdfauthor=, pdfsubject=, pdfkeywords=}
\usepackage{tabularx}
\usepackage[table]{xcolor} % Table colors
\usepackage{scrextend}
\addtokomafont{labelinglabel}{\sffamily}
\usepackage{listings}
\usepackage{adjustbox}
\usepackage{color}

% Define colors
\definecolor{ashgrey}{rgb}{0.7, 0.75, 0.71}

% Listing style
\lstset{
  backgroundcolor=\color{ashgrey}, % choose the background color; you must add \usepackage{color} or \usepackage{xcolor}; should come as last argument
  basicstyle=\footnotesize,        % the size of the fonts that are used for the code
  breakatwhitespace=true,          % sets if automatic breaks should only happen at whitespace
  breaklines=true,                 % sets automatic line breaking
  extendedchars=true,              % lets you use non-ASCII characters; for 8-bits encodings only, does not work with UTF-8
  frame=single,	                   % adds a frame around the code
  rulecolor=\color{black},         % if not set, the frame-color may be changed on line-breaks within not-black text (e.g. comments (green here))
  keepspaces=true,                 % keeps spaces in text, useful for keeping indentation of code (possibly needs columns=flexible)
  columns=fullflexible,		   % make copy and paste possible
  showstringspaces=false,          % if true show spaces in strings adding particular underscores
  showspaces=false,                % if true show spaces everywhere adding particular underscores; it does not override 'showstringspaces'
}

% Uncomment for production
% \syntaxonly

% Style
\pagestyle{headings}

% Turn on indexing
\makeindex[intoc]

% Define document
\author{D. Leeuw}
\title{TCP: DHCP}
\date{\today\\v.0.8.0}

\begin{document}
\selectlanguage{dutch}

\maketitle

\copyright\ 2023 Dennis Leeuw\\

\begin{figure}
\includegraphics[width=0.3\textwidth]{CC-BY-SA-NC.png}
\end{figure}

\bigskip

Dit werk is uitgegeven onder de Creative Commons BY-NC-SA Licentie en laat anderen toe het werk te kopi\"eren, distribueren, vertonen, op te voeren, en om afgeleid materiaal te maken, zolang de auteurs en uitgever worden vermeld als maker van het werk, het werk niet commercieel gebruikt wordt en afgeleide werken onder identieke voorwaarden worden verspreid.


%%%%%%%%%%%%%%%%%%%
%%% Introductie %%%
%%%%%%%%%%%%%%%%%%%

\frontmatter
\chapter{Over dit Document}
\section{Leerdoelen}
Dit document gaat over ethernet. Het leert je wat ethernet is. Om meer te weten te komen over wat er op je netwerk gebeurt maak je ook kennis met de tools: \texttt{arp} en \texttt{wireshark}.

Na het bestuderen van dit document heb je een basis kennis van ethernet, weet je wat het adres-formaat is van ethernet en kan je informatie over het netwerk opvragen via verschillende tools.

In de praktijk opdrachten leer je om je eigen kabel te maken en deze te testen. Tevens maak je je eigen Wi-Fi netwerk.

\section{Voorkennis}
Er wordt van de lezer verwacht dat deze kennis heeft van de volgende onderwerpen:
\begin{itemize}
\item TCP/IP
\item CMOS en BIOS tijd
\end{itemize}


%%%%%%%%%%%%%%%%%
%%% De inhoud %%%
%%%%%%%%%%%%%%%%%
\tableofcontents

\mainmatter
\chapter{Inleiding}
TCP/IP is het protocol van het Internet en tegenwoordig ook van de meeste andere computer-netwerken. Vroeger waren er veel verschillende protocollen. Van die verschillende protocolen is eigenlijk alleen TCP/IP overgebleven.

TCP/IP is ontworpen als een protocol dat moest werken over onbetrouwbare netwerkverbindingen en dat ook bij een kernwapen-oorlog waarbij delen van het netwerk uitvallen toch nog zijn data moet kunnen versturen van A naar B. Het moet dus feitelijk zelf uitzoeken hoe data het beste van A naar B kan komen.

\section{Het DHC-Protocol}
DHCP maakt gebruik van broadcasts tussen de client en server om elkaar te kunnen vinden. RFC 2131 beschrijft DHCP, wij zullen het hier doen met een minimale beschrijving. DHCP gebruikt een UDP op de transport-laag.

Een werkstation dat geconfigureerd is voor automatische IP toewijzing (DHCP) zal tijdens het configureren van het netwerk (opstart fase) een IP-broadcast (255.255.255.255) het netwerk op sturen, dit is de zogenaamde DHCP Discover. Dit bericht wordt door alle machines op het netwerk gehoord en is gericht aan port 67, afkomstig van port 68.

Elke DHCP-server op het netwerk zal een antwoord terug sturen, DHCP Offer. Dit bericht mag een unicast of een broadcast zijn.

De client ontvangt de packetten in een bepaalde volgorde en er wordt aangenomen dat het eerste packet dat binnenkomt van de dichtstbijzijnde DHCP server is. Dus de client zal proberen om het daarin aangeboden IP-adres te claimen voor gebruik. Het doet dit door een DHCP Request uit te sturen als broadcast om bevestiging te krijgen dat het dit IP-adres mag gebruiken. Alle andere DHCP-servers weten nu dat hun aangeboden IP-adres niet gebruikt gaat worden en dat ze deze dus weer in de pool van te gebruiken IP-adressen kunnen zetten.

De DHCP-server die hoort dat er voor zijn offer nu een request binnen komt, bevestigd naar de client dat het okee is om dit IP-adres te gebruiken. Dit bericht mag weer een unicast of een broadcast zijn.

De berichten van de client naar de server moeten broadcasts zijn zodat alle DHCP servers weten wat er gebeurd, van de server naar de client mogen het ook unicast berichten zijn om in complexe netwerken de hoeveelheid verkeer te minimaliseren.


\section{Lease time}
Een werkstation krijgt een IP-adres voor een bepaalde tijd (lease time). Een DHCP-server heeft een beperkte set van IP-adressen die uitgedeeld mogen worden. Door een lease-time in te stellen weet de DHCP-server dat zelfs als een werkstation van het netwerk verdwijnt dat het na een bepaalde tijd het IP-adres kan hergebruiken. Een client die een IP-adres wil blijven gebruiken moet aan de server vragen of dit mag.


\section{DHCP Options}
Met DHCP kan je meer uitdelen dan alleen IP-adressen. Standaard wordt er en IP-adres, een subnet-mask en een default gateway uitgedeeld. In de meeste netwerken wordt ook aangegeven wat de IP-adressen van de name-servers zijn.

Het is ook toegestaan om bijvoorbeeld de hostname, de domainname of extra routes mee te geven.

Totslot zijn er nog opties die het mogelijk te maken om de client te vertellen welke NTP-server of welke mail-server er gebruikt kan worden. Er zijn nog veel meer opties die uitgedeeld kunnen worden, maar deze laatste worden maar weinig gebruikt in de huidige netwerken.



\chapter{Scope}
Binnen een netwerk is het handig als je bepaalde opties op je DHCP-server kan groeperen, zodat je bijvoorbeeld naar iedereen dezelfde default gateway stuurt en dat niet voor elke host op je netwerk hoeft aan te maken. Zo'n groepering heet een scope. Scope kan het best vertaald worden met bereik of omvang.

\section{Range}
De range die een DHCP-server heeft zijn de IP-adressen die een server uit mag delen. Meestal is dit een reeks op elkaar volgende IP-adressen. Bijvoorbeeld 192.168.1.100 tot 192.168.1.200. 

\section{Option: subnet-mask}
Bij een reeks hoort natuurlijk een subnet-mask, in het geval van een klasse-C adres is dat 255.255.255.0.

\section{Option: routers}
De meeste netwerken hebben een router die zorgt voor de verbinding met andere netwerken. Deze router staat bekend als de default gwateway. Ook deze optie moet gezet worden in de scope zodat voor de netwerk range bekend is wie er voor de verbinding naar buiten zorgt.

\section{Option: domain-name}
De client wil ook graag weten in welk domain het leeft. De domain-name option is ervoor om dat aan de client duidelijk te maken.

\section{Option: domain-name-servers}
Om domeinnamen te vertalen naar IP-adressen moet de client weten welke domain-name-servers er zijn op het netwerk. Die de optie domain-name-servers worden \'e\'en of meerdere servers opgegeven.



\chapter{Een dynamisch vast IP-adres}
Het is via DHCP ook mogelijk om een bepaalde computer altijd hetzelfde IP-adres te geven. In DHCP vertellen we de server dan dat aan een bepaald MAC-adres altijd hetzelfde IP-adres gegeven moet worden. Er wordt dus een reservering gemaakt voor een MAC-adres. Als dit MAC-adres om een IP-adres vraagt krijgt het het gereserveerde adres.

Zo'n configuratie voor een vast IP adres behorende bij een MAC-adres wordt ook wel een host declaratie genoemd.

\section{Option host-name}
Binnen een host declaratie kan je ook de host-name optie opgeven, zodat de machine naast een vast IP-adres ook een vaste host-name heeft.


\chapter{Opdrachten}
\section{DHCP-server op Windows}
Installeer een DHCP-server op Windows door gebruik te maken van het document op \url{https://learn.microsoft.com/nl-nl/troubleshoot/windows-server/networking/install-configure-dhcp-server-workgroup}. Dit kan op een fysieke server of op een virtual machine.


%%%%%%%%%%%%%%%%%%%%%
%%% Index and End %%%
%%%%%%%%%%%%%%%%%%%%%
\backmatter
\printindex
\end{document}

%%% Last line %%%
